\begin{abstract}
新能源发电的波动性给电力系统运行带来新的挑战。本文以含新能源和储能的电力系统为研究对象，建立考虑功率平衡、设备容量与运行成本的优化模型，分析新能源出力变化对系统调度的影响。

研究采用数学建模与数值分析相结合的方法，通过统一描述发电、负荷和储能设备的运行约束，构建具有可解释性的优化框架。模型以系统运行成本最小为目标，并对储能充放电行为进行约束，以保证调度方案满足设备运行条件。

本文给出模型的主要变量、目标函数和约束关系，并讨论求解结果的物理意义。研究表明，通过协调新能源发电与储能资源，可以改善电力系统的运行方式。所建立的框架为后续开展不确定性分析和多时间尺度协调提供基础。

\end{abstract}
\begin{abstractEn}
Renewable generation introduces new challenges to power system operation. This study develops an optimization model for a power system with renewable generation and energy storage, subject to power balance, capacity limits, and operating constraints.

The model combines mathematical modeling with numerical analysis. Its objective is to minimize system operating costs while coordinating generation, demand, and storage. The physical meaning of the variables and constraints is discussed to support interpretation of the scheduling results.

The main variables, objective function, and constraints are presented, and their physical meaning is discussed. Coordinating renewable generation and energy storage provides a framework for subsequent uncertainty analysis and coordination across multiple time scales.
\end{abstractEn}
